\section{Methods}
\label{sec:methods}

\subsection{Overview and preprocessing}
\label{subsec:preproc}

All models predict HS at 1\,m resolution from the predictor stack described in
Section~\ref{subsec:channels}. Each channel is normalised to a comparable
range using fixed, physically motivated scalings (e.g.\ elevation is
standardised with a reference mean and standard deviation, slopes are divided
by $90^\circ$, the TPI is clipped to $[-1,1]$, the $S_x$ indices are divided by
$90^\circ$, persistence channels are already in $[0,1]$, and SCE is binarised).
Invalid sentinel values ($-9999$) and NaNs are replaced by zero after
normalisation. Pixels without valid HS are masked out of the loss and of every
metric. Crucially, normalisation is applied to the full 22-channel stack
\emph{before} any channel subset is selected, so that the ablation experiments
(Section~\ref{subsec:res_ablation}) do not alter the normalisation statistics
of the remaining channels.

\subsection{Models}
\label{subsec:models}

\paragraph{Random Forest (baseline).}
As a strong non-spatial baseline we train a Random Forest regressor
\citep{breiman2001} on the per-pixel feature vectors. Hyperparameters were
optimised with Optuna \citep{akiba2019} and the final configuration uses
300 trees, a maximum depth of 15, a minimum of one sample per leaf,
\texttt{log2} feature subsampling and a minimum of two samples to split. The
RF is trained on the union of the training and validation pixels (standard
refit after hyperparameter optimisation) and predicts each pixel
independently; spatial metrics are computed by reassembling the per-pixel
predictions into their original tiles.

\paragraph{U-Net.}
The U-Net \citep{ronneberger2015} is a symmetric encoder--decoder with skip
connections. Our implementation uses a base width of 48 feature maps and the
standard four-level contracting/expanding path, and is trained end-to-end as a
pixel-wise regressor. We use Group Normalisation in the U-Net as well, since
Batch Normalisation proved unreliable under our small batch sizes -- its running
statistics collapsed at evaluation time -- and using the same normalisation in
both networks ensures that the U-Net/ResUNet++ comparison isolates the effect of
the architecture rather than of the normalisation layer.

\paragraph{ResUNet++.}
ResUNet++ \citep{jha2019} augments the residual U-Net
\citep{zhang2018} with squeeze-and-excitation blocks in the encoder, an atrous
spatial pyramid pooling (ASPP) bridge, and attention gates in the decoder.
These components improve gradient flow and multi-scale context aggregation. We
use a base width of 64 and Group Normalisation (8 groups) instead of Batch
Normalisation, which we found essential for stable training under the small
batch sizes imposed by high-resolution tiles and the heterogeneous,
date-dependent statistics of the snowpack.

\begin{figure}[t]
  \centering
  \includegraphics[width=\linewidth]{fig_architecture.pdf}
  \caption{ResUNet++ adapted for pixel-wise snow-depth regression. A three-level
  residual encoder with squeeze-and-excitation (SE) blocks and Group
  Normalisation encodes the $N$-channel input; an atrous spatial pyramid pooling
  (ASPP, dilations $r=1,6,12,18$) bottleneck aggregates multi-scale context; and
  a decoder with attention-gated skip connections (AG) reconstructs the field.
  The $1\times1$ output head is linear (no activation), yielding a continuous
  snow-depth map. Resolutions and channel counts correspond to the
  base-width-64 configuration on $256\times256$ tiles.}
  \label{fig:architecture}
\end{figure}

\subsection{Adapting ResUNet++ for pixel-wise snow-depth regression}
\label{subsec:adaptation}

ResUNet++ was originally proposed for binary medical image
\emph{segmentation} \citep{jha2019}; running it as a snow-depth predictor
required adapting it to the pixel-wise regression framework introduced by
MAPunet \citep{betato2025mapunet}. Concretely, the model learns a mapping
$f:\mathbb{R}^{H\times W\times N}\!\rightarrow\!\mathbb{R}_{\ge 0}^{H\times W}$
that turns an $N$-channel input volume (here $N=22$, or $26$ with the
meteorological channels) into a single continuous snow-depth field
(Fig.~\ref{fig:architecture}). We made the following modifications:

\begin{itemize}
  \item \textbf{Input layer.} The first convolution was widened to accept the
  $N$ co-registered predictor channels (Table~\ref{tab:channels}) instead of a
  fixed three-band image, so that arbitrary channel subsets can be fed to the
  network (used by the ablation in Section~\ref{subsec:res_ablation}).

  \item \textbf{Regression output.} The final $1\times1$ convolution produces a
  single feature map and, crucially, carries \emph{no} output activation (no
  sigmoid/softmax), yielding unbounded continuous values; predictions are
  clipped to non-negative depths only at inference.

  \item \textbf{Masked regression loss.} High-resolution LiDAR maps contain
  \textit{nodata} pixels (occlusions, gaps, out-of-catchment cells). Following
  the masked-supervision principle of MAPunet, the loss and all metrics are
  evaluated only over valid pixels; invalid pixels are excluded from both the
  pixel-wise MSE and the spatial-correlation term of
  Eq.~\eqref{eq:loss}.

  \item \textbf{Normalisation for heterogeneous inputs.} Group Normalisation
  (8 groups) replaces Batch Normalisation to remain stable under the small
  batches imposed by high-resolution tiles and the strongly date-dependent
  statistics of the snowpack, and every channel is scaled with fixed physical
  factors (Section~\ref{subsec:preproc}) before being fed to the network.
\end{itemize}

\noindent The same adaptations apply to the U-Net baseline, so that the two
architectures differ only in their internal blocks and are compared on an equal
footing.

\subsection{Spatial evaluation metrics}
\label{subsec:metrics}

\paragraph{Pixel-wise metrics.}
We report the coefficient of determination (\Rsq{}, equivalently the
Nash--Sutcliffe efficiency), the mean absolute error (MAE), the root mean
squared error (RMSE) and the bias, all computed over the valid test pixels.

\paragraph{SPAEF.}
To assess the fidelity of the predicted \emph{spatial pattern} we use the
Spatial Efficiency metric \citep{koch2018}, defined as
%
\begin{equation}
  \SPAEF = 1 - \sqrt{(\alpha-1)^2 + (\beta-1)^2 + (\gamma-1)^2},
  \label{eq:spaef}
\end{equation}
%
where the three components measure complementary aspects of pattern agreement:
%
\begin{equation}
  \alpha = \rho(o,s), \qquad
  \beta  = \frac{\sigma_s/\mu_s}{\sigma_o/\mu_o}, \qquad
  \gamma = \frac{\sum_{j=1}^{K}\min\!\big(h^{o}_{j}, h^{s}_{j}\big)}
                {\sum_{j=1}^{K} h^{o}_{j}} .
  \label{eq:spaef_components}
\end{equation}
%
Here $o$ and $s$ are the observed and simulated fields, $\rho$ is the Pearson
correlation coefficient ($\alpha$ measures linear co-location of the pattern),
$\beta$ is the ratio of the coefficients of variation (it penalises
over-/under-dispersed patterns), and $\gamma$ is the intersection of the
$K$-bin histograms $h^{o}$ and $h^{s}$ of the two fields (it measures the
agreement of the value distributions). A perfect match yields
$\SPAEF=1$; the metric is unbounded below. \SPAEF{} is computed per tile and
averaged over the test set.

\paragraph{MSPAEF.}
As a second pattern-oriented metric we use the modified \SPAEF{} (\MSPAEF{}) of
\citet{karpasitis2026},
%
\begin{equation}
  \MSPAEF = 1 - \tfrac{1}{4}\big[(\alpha-1)^2 + \beta^2 + \gamma^2 + \delta^2\big],
  \label{eq:mspaef}
\end{equation}
%
where $\alpha=\rho(o,s)$ is again the Pearson correlation, $\beta$ is the RMSE
normalised by the interquartile range (IQR) of the observations, $\gamma$ is
the absolute mean bias normalised by the same IQR, and $\delta$ is a
standard-deviation-ratio term. Compared with the classical \SPAEF{} of
Eq.~\eqref{eq:spaef}, \MSPAEF{} (i) requires no histogram-bin parameter and
(ii) is explicitly sensitive to the absolute mean bias, which \SPAEF{} ignores.
This is particularly relevant here, since all models systematically
under-predict the anomalously snowy 2025 test year; we therefore report
\MSPAEF{} as a complementary, bias-aware indicator. Like \SPAEF{}, it is
computed per tile and averaged over the test set.

\subsection{Hybrid spatial loss}
\label{subsec:loss}

Training a CNN with a pure pixel-wise MSE encourages regression to the mean and
tends to produce over-smoothed maps. To explicitly reward spatial pattern
fidelity, we train the CNNs with a hybrid loss that linearly combines the
pixel-wise MSE with a spatial Pearson-correlation term:
%
\begin{equation}
  \mathcal{L} = (1-\lpear)\,\underbrace{\mathrm{MSE}(o,s)}_{\text{pixel-wise}}
  + \lpear\,\underbrace{\big(1-\rho_{\text{spatial}}(o,s)\big)}
  _{\text{pattern}},
  \label{eq:loss}
\end{equation}
%
where $\rho_{\text{spatial}}$ is the Pearson correlation between the predicted
and observed fields computed over the valid pixels of each tile, and
$\lpear\in[0,1]$ controls the trade-off between absolute accuracy and pattern
agreement. With $\lpear=0$ the loss reduces to plain MSE; increasing $\lpear$
progressively prioritises the reproduction of the spatial pattern. The
correlation term is differentiable and adds negligible computational cost. We
characterise the effect of $\lpear$ through a full sweep
(Section~\ref{subsec:res_lambda}).

\subsection{Training and evaluation protocol}
\label{subsec:protocol}

\paragraph{Implementation and hardware.}
All models were implemented in PyTorch~2.11 (CUDA~12.8) and trained on a single
desktop workstation with an NVIDIA GeForce RTX~5060~Ti GPU (16\,GB), an AMD
Ryzen~5 3400G CPU and 32\,GB of RAM, running Windows~10. The Random Forest used
scikit-learn~1.8, and all hyperparameter searches used Optuna~4.8
\citep{akiba2019}.

\paragraph{Hyperparameter optimisation.}
All model hyperparameters were tuned with Optuna \citep{akiba2019} using a
Tree-structured Parzen Estimator (TPE) sampler on the train/validation split,
keeping the 2025 test year strictly untouched during the search. For the CNNs we
jointly searched the base width, optimiser, learning rate, weight decay,
dropout, batch size and gradient clipping (Table~\ref{tab:hyperparams}); each
trial was trained for up to 150 epochs with a median pruner (10 start-up trials,
20 warm-up epochs) and scored by the validation MSE. For the Random Forest we
searched the number of trees, the maximum depth, the minimum number of samples
per leaf and per split, and the feature-subsampling rule, scoring trials by the
validation \Rsq{}. The best configuration of each model (Table~\ref{tab:hyperparams})
was then fixed for all subsequent experiments; the spatial-loss weight $\lpear$
was tuned separately (Section~\ref{subsec:res_lambda}) on top of the chosen
ResUNet++ backbone, and the final models were retrained for 50 epochs (see
below).

\begin{table}[t]
  \centering
  \caption{Hyperparameter search spaces (Optuna, TPE sampler) and selected
  configurations. The CNN search minimised the validation MSE over up to 150
  epochs with median pruning; the Random Forest search maximised the validation
  \Rsq{}.}
  \label{tab:hyperparams}
  \small
  \begin{tabular}{lccc}
    \toprule
    \multicolumn{4}{l}{\textit{CNNs (U-Net and ResUNet++)}}\\
    Hyperparameter & Search space & U-Net & ResUNet++ \\
    \midrule
    Base width        & $\{32, 48, 64\}$                      & $48$ & $64$ \\
    Optimiser         & $\{$Adam, AdamW$\}$                   & Adam & AdamW \\
    Learning rate     & $[3{\times}10^{-5},\,5{\times}10^{-4}]$ (log) & $3.06{\times}10^{-5}$ & $1.29{\times}10^{-4}$ \\
    Weight decay      & $[5{\times}10^{-6},\,6{\times}10^{-4}]$ (log) & $2.15{\times}10^{-5}$ & $1.24{\times}10^{-4}$ \\
    Dropout           & $[0,\,0.20]$                          & $0.001$ & $0.077$ \\
    Batch size        & $\{8, 16\}$                           & $16$ & $8$ \\
    Gradient clipping & $\{0,\,1.0\}$                         & $1.0$ & $1.0$ \\
    \midrule
    \multicolumn{2}{l}{\textit{Random Forest}} & \multicolumn{2}{c}{Selected} \\
    \midrule
    No.\ of trees      & $\{100, 200, 300, 500\}$            & \multicolumn{2}{c}{$300$} \\
    Max.\ depth        & $\{10, 15, 20, 30, \text{None}\}$   & \multicolumn{2}{c}{$15$} \\
    Min.\ samples leaf & $\{1, 5, 10, 20\}$                  & \multicolumn{2}{c}{$1$} \\
    Max.\ features     & $\{$sqrt, log2, $0.3$, $0.5\}$      & \multicolumn{2}{c}{log2} \\
    Min.\ samples split& $\{2, 5, 10\}$                      & \multicolumn{2}{c}{$2$} \\
    \bottomrule
  \end{tabular}
\end{table}

\paragraph{Training.}
All CNN models are trained for 50 epochs without early stopping and without
data augmentation, to keep the comparison controlled and reproducible. Every
model is evaluated at the checkpoint with the lowest validation loss
(best-validation epoch), applying the same selection rule to all architectures
for a fair comparison.

\paragraph{Reproducibility and seeds.}
Every CNN experiment is repeated with three random seeds (42, 123, 7), with the
seed fixed immediately before model construction. We report the mean and
standard deviation across seeds for all metrics; the seed spread is an
essential part of our results, as it reveals stability differences that
single-run numbers would mask. The RF, being almost deterministic for this
problem, is likewise evaluated over the three seeds for consistency.

\paragraph{Experiments.}
Four experiments are reported in Section~\ref{sec:results}:
(i) a head-to-head comparison of RF, U-Net and ResUNet++ on the 22-channel
dataset;
(ii) a sweep of the spatial-loss weight $\lpear\in\{0,0.1,0.25,0.4,0.5,0.6,
0.75,1.0\}$ for ResUNet++ (eight values $\times$ three seeds);
(iii) a leave-one-group-out ablation of the predictor families (six
configurations $\times$ three seeds); and
(iv) an analysis of the 26-channel meteorological configuration.
